\section{Discussion and lessons learned}
\label{sec:discussion}
The catalogue records where a document should be found, the index determines
whether search can retrieve it, and the evidence pack shows which text
actually reached the agent. Agreement at one stage does not establish
success at the next.

HMRS combines workflow-mediated, content-informed collection placement and
document descriptions with query-time scope navigation and index probes after
changes. The results do not show that this arrangement outperforms a flat,
faceted index. They show how the shared workflow makes it possible to inspect
collection choice, index visibility, text coverage and answer evaluation at
different stages. A comparison with a flat alternative using the same metadata,
rankers and source-reading tools is needed to attribute any benefit specifically
to the hierarchy.

\textbf{Check visibility after a write.} The rename-and-move probe found
stale index entries. Comparing a retrieval probe with the index helps
locate the fault; a stale response does not by itself establish an index
problem. Neither check replaces reconciliation of the files, tree registry and
library metadata \citep{park-metadata-quality-2009,paivarinta-ecm-2005}.

\textbf{Check the text as well as the route.} A0 reached the labelled
collection on every question but missed many target strings. Under the same route,
A0h contained more target strings than A0, but this is a configuration-level
text-coverage result. It does not identify the separate contributions of the
extra vector search, scored document text or prioritised windows. Nor does the
comparison with PD isolate an effect of hierarchy. In the fixed-pool replay,
verifier-score ranking displaced gold documents. Passage admission and
document order therefore need separate tests.

The native answer-quality comparison is descriptive. It did not exercise
catalogue-first routing, used non-matched search scopes and source access, and
did not yield a corrected significant paired difference. Likewise, operational
lesson extraction, citation markers and instructed abstention were only
partially observed. They show why scope and completion status should be
reported, not that the server guarantees supported answers.

\textbf{Report the actual scope.} A scope report should distinguish requested
scope from effective scope and record skipped collections, transport
truncation or pagination, and source-read or scoring failures. The E7
search returned success after skipping eight unreadable collections. In the
other direction, the E2 agent sometimes searched beyond its configured
collections and could follow up with source reads. Those capabilities are part
of the deployed answering workflow, not an advantage attributable to the
ranker. A scope report lets readers distinguish a completed search without
evidence from one that never reached part of the corpus.
